% ==============================================================================
% SEZIONE 3: KPI MACROECONOMICI E INDICATORI DI SCENARIO
% ==============================================================================

\section{Market Drivers}

\begin{frame}{KPI Macroeconomici e Indicatori di Scenario}
    Le variabili esogene definiscono il perimetro d'azione delle utility, legando la profittabilità al differenziale tra prezzi spot e domanda elettrica nelle aree geografiche Tier 1.
    \vspace{0.1cm}
    
    \note{I market drivers evidenziano uno scenario macroeconomico asimmetrico e a due velocità tra i vari paesi core.}
    
    \begin{table}
        \centering
        \small
        \begin{tabular}{lcccc}
            \toprule
            \textbf{Area Tier 1} & \textbf{GDP (\%)} & \textbf{CPI (\%)} & \textbf{Spot Price (€/MWh)} & \textbf{Demand (TWh)} \\
            \midrule
            \textbf{Italia}  & 0.65 & 1.63 & 115.3 & 311.3 \\
            \textbf{Iberia}  & 2.80 & 2.69 & 65.5  & 320.7 \\
            \textbf{Brasile} & 2.28 & 5.02 & 31.9  & 754.4 \\
            \textbf{Cile}    & 2.40 & 4.21 & 53.3  & 85.1  \\
            \textbf{USA}     & 2.20 & 2.72 & --    & 4503.8 \\
            \bottomrule
        \end{tabular}
        \caption{Key Market Drivers e asimmetrie regionali nel FY 2025.}
    \end{table}
    \vspace{-0.1cm}
    \pause
    
    \begin{alertblock}{Esposizione al Rischio di Cambio (FX)}
        L'esposizione al rischio di cambio per il Gruppo rimane significativa, con \textbf{€12,7 miliardi esposti}: il monitoraggio del CPI e dei tassi locali è un fattore critico di gestione.
    \end{alertblock}
    \note{L'esposizione al rischio FX di dodici virgola sette miliardi impone un monitoraggio costante di CPI e tassi locali.}
\end{frame}


\begin{frame}{Valutazione delle Divergenze Strategiche}
    L'analisi degli indicatori di scenario evidenzia una marcata asimmetria competitiva tra le diverse geografie core:
    \vspace{0.3cm}
    
    \note{L'analisi di scenario mostra tre andamenti macroeconomici distinti.}
    
    \begin{itemize}
        \item \textbf{Iberia:} Scenario ottimale con crescita della domanda robusta (320,7 TWh) accoppiata a prezzi spot competitivi (€65,5/MWh).
        \pause
        \note{L'Iberia mostra uno scenario eccellente con forte incremento dei consumi e prezzi all'ingrosso competitivi.}
        
        \item \textbf{Italia:} Mantiene prezzi strutturalmente elevati (€115,3/MWh) a fronte di una domanda stagnante influenzata da un PIL contenuto (0.65\%).
        \pause
        \note{L'Italia presenta invece margini sotto pressione, con prezzi spot elevati e consumi frenati dalla crescita debole.}
        
        \item \textbf{USA:} Emergono come il mercato di riferimento per scala dimensionale (4.503,8 TWh), giustificando la forte concentrazione degli investimenti.
    \end{itemize}
    \vspace{0.2cm}
    \pause
    \note{Gli Stati Uniti si confermano il mercato chiave per volumi energetici, accentrando la quota maggiore di investimenti.}
    
    \begin{block}{Hedge Naturale}
        La diversificazione geografica globale funge da \textbf{protezione (hedge) naturale} contro la volatilità dei singoli mercati e regimi normativi nazionali.
    \end{block}
    \note{Questa bilanciatura geografica funge da hedge naturale, stabilizzando i flussi finanziari aggregati contro gli shock locali.}
\end{frame}
